\section{导读：为什么这期是 AI 应用创业的路线笔记}

本节先定位这期。EP95 是一场“接力式访谈”：第一次发生在 2024 年秋天，肖弘刚完成融资，主要讨论连续创业、资本、Monica 和 AI 应用方法论；第二次发生在 2025 年春节后，DeepSeek 改写中国 AI 应用底层生态，肖弘开始讲即将发布的 Manus，也就是国内 Agent 产品的“第一枪”。因此这期不是一段静态创始人故事，而是一份 AI 应用创业者在模型能力急剧变化中的连续思考记录。

\begin{importantbox}{本期核心命题}
肖弘的主线可以压缩成三句话：第一，AI 应用公司要预判模型能力和大厂动作，提前把产品容器准备好；第二，模型能力会技术平权，但应用公司的壁垒来自场景、品牌、用户理解、反馈和执行；第三，世界不是线性外推，Founder 要用博弈方式思考，让自己成为局面中的重要变量。
\end{importantbox}

\begin{knowledgebox}{视觉策略说明}
本视频是固定访谈画面，没有投屏或产品演示。正文只使用封面作为来源识别；正文图像全部为自制概念图，用来解释接力访谈结构、VC 成本、预判大厂、Monica、AI 应用分类、Manus Agent 产品和 Founder 博弈。
\end{knowledgebox}

\begin{figure}[H]
\centering
\includegraphics[width=0.96\textwidth]{xiaohong-relay-interview.png}
\caption{肖弘接力式访谈地图：两次访谈记录 AI 应用创业者在不稳定环境中的连续思考。自制概念图，依据全片章节结构整理。}
\end{figure}

\begin{knowledgebox}{读图：这不是一条直线，而是两次状态截面}
图中从 2024 秋融资后，到第一段创业、Monica、2025 春节 DeepSeek 冲击，再到 Manus。读这张图时要注意：同一个 Founder 在不同时间点面对的外部模型能力、资本环境、产品机会和组织问题都不同，因此判断也在变化。
\end{knowledgebox}

\subsection{阅读路线}

本节给出阅读路线。整期可以拆成六个问题：第一，肖弘第一段创业给他留下了哪些产品、商业化和资本教训？第二，为什么“VC 是很贵的集资手段”？第三，Monica 为什么选择浏览器插件、海外市场和 all-in-one AI assistant？第四，AI 应用如何分类：主场景补充、模型能力变化、垂直领域外溢？第五，Manus 为什么被理解为国内 Agent 第一枪？第六，Founder 为什么要从逻辑推理转向博弈思维？

\noindent\begin{tabularx}{\textwidth}{p{0.20\textwidth}p{0.34\textwidth}X}
\toprule
阅读问题 & 访谈材料 & 要形成的判断 \\
\midrule
第一段创业学到什么 & 微信公众号工具、SCRM、商业化、出售公司 & 工具型产品要早做商业化，别被大故事带偏。 \\
资本如何影响决策 & VC 很贵、投资人权重、资本维度 & 融资是工具，不是老板，也不是战略本身。 \\
Monica 为什么成立 & 插件、海外、ChatGPT for Google、用户反馈 & AI 应用要卡住模型能力外溢后的产品容器。 \\
应用公司怎么找机会 & Jasper、ChatGPT、Cursor、DeepSeek、Manus & 观察模型能力和用户主场景之间的缝隙。 \\
Manus 的产品意义 & 任务、工具、环境、记忆、交付结果 & Agent 从聊天走向能做事的系统。 \\
Founder 怎么思考 & 非线性、博弈、时代年龄、贪嗔痴 & 创业不是线性推理，要改变局面。 \\
\bottomrule
\end{tabularx}

\begin{warningbox}{时间边界}
第一次访谈发生在 2024 年秋天。那部分说“今年”多指 2024 年，说“去年”多指 2023 年。第二次访谈发生在 2025 年春节后，讨论的是 DeepSeek 爆发后、Manus 上线前的状态。
\end{warningbox}

\subsection{本章小结}

EP95 是一份 AI 应用创业方法论笔记。它和 EP128 Peak/Manus 最后访谈、EP101 YouWare、EP110 Agent 技术报告可以连读：前者讲 Manus 后来的组织和收购前状态，这期讲 Manus 发布前的创始人思考和应用创业底层方法。

